
% Funciones propias
\newcommand{\testcase}[3]{
	\begingroup
	\begin{table}[H]
		\centering
		\caption{#1} % Parámetro 1: El título de la tabla
		\label{#2}   % Parámetro 2: El identificador para referenciar (\ref{})
		\vspace{6pt}
		\renewcommand{\arraystretch}{1.4} 
		
		\begin{tabularx}{\textwidth}{ | >{\bfseries}p{4.5cm} | X | }
			\hline
			#3 % Parámetro 3: Todo el contenido interno de la tabla
			\hline
		\end{tabularx}
	\end{table}
	\endgroup
}


% Función para generar tablas de Descripción de Casos de Uso
% Parámetro 1: Nombre del Caso de Uso (para el título)
% Parámetro 2: Label para referenciar (\ref{})
% Parámetro 3: Contenido de la tabla
\newcommand{\usecase}[3]{
	\begingroup
	\begin{table}[H]
		\centering
		\caption{#1} 
		\label{#2}   
		\renewcommand{\arraystretch}{1.3} 
		
		% m{3.5cm} -> Ancho fijo para los nombres de los campos
		% X -> El resto del ancho de la página para la descripción
		% >{\bfseries} -> Pone en negrita automático la primera columna
		\begin{tabularx}{\textwidth}{ | >{\bfseries}m{3.5cm} | X | }
			\hline
			#3 
			\hline
		\end{tabularx}
	\end{table}
	\endgroup
}

\NewDocumentCommand{\tablaCasoPrueba}{m m m m +m}{%
	% ---- Bloque de encabezado ----
	\noindent
	\begin{tblr}{
			width   = \linewidth,
			colspec = {Q[l,font=\bfseries,wd=3.8cm] X[l]},
			hlines, vlines,
			rowsep  = 3pt,
		}
		\SetCell[c=2]{l} \textbf{Caso de Prueba} & \\
		Nombre                  & #1 \\
		Objetivo de la Prueba   & #2 \\
		Descripción de la prueba & #3 \\
		Condiciones de entrada  & #4 \\
		\SetCell[c=2]{l} \textbf{Combinaciones de valores} & \\
	\end{tblr}
	
	% ---- Bloque de combinaciones ----
	\noindent
	\begin{tblr}{
			width   = \linewidth,
			colspec = {Q[c,wd=0.6cm] X[2,l] X[1,l] X[3,l] X[2,l] X[2,l]},
			hlines, vlines,
			rowsep  = 3pt,
		}
		\textbf{CP} &
		\textbf{Escenarios} &
		\textbf{Variables} &
		\textbf{Valores de entrada} &
		\textbf{Resultado esperado} &
		\textbf{Resultado real} \\
		#5
	\end{tblr}
}

\ExplSyntaxOn

\newcounter{reqcounter}
\setcounter{reqcounter}{0}
\renewcommand{\thereqcounter}{RF-\arabic{reqcounter}}
\renewcommand\theadfont{\bfseries} % para que \thead use negrita

\tl_new:N \g_req_rows_tl

\NewDocumentCommand{\requisitos}{ m m }
{
	\subsubsection*{#1}
	
	% 1. Limpiar acumulador
	\tl_gclear:N \g_req_rows_tl
	
	% 2. Partir la lista y acumular filas
	\seq_set_split:Nnn \l_tmpa_seq { ; } { #2 }
	\seq_map_inline:Nn \l_tmpa_seq
	{
		\tl_set:Nx \l_tmpb_tl { \tl_trim_spaces:n { ##1 } }
		\tl_if_empty:NF \l_tmpb_tl
		{
			% Guardamos el texto del requisito expandido (esto SÍ se expande aquí)
			\tl_gput_right:Nx \g_req_rows_tl
			{
				\exp_not:n { \hline \refstepcounter{reqcounter} }
				% \thereqcounter y \label se expanden en tiempo de volcado (dentro de la tabla)
				\exp_not:n { \thereqcounter }
				\exp_not:n { \label{rf:\arabic{reqcounter}} }
				\exp_not:n { & }
				\exp_not:V \l_tmpb_tl
				\exp_not:n { \\ }
			}
		}
	}
	
	% 3. Volcar dentro de la tabla
	\begin{longtable}{|p{0.12\linewidth}|p{0.82\linewidth}|}
		\caption{ #1 }
		\label{tab:frameworks-backend} \\
		
		\hline
		\textbf{ID} & \thead[l]{Descripción} \\
		\hline
		\endfirsthead
		\multicolumn{2}{c}{\tablename\ -- continuación} \\
		\hline
		\textbf{ID} & \thead[l]{Descripción} \\
		\hline
		\endhead
		\hline
		\endfoot
		\hline
		\endlastfoot
		\g_req_rows_tl
		\hline
	\end{longtable}
}

\ExplSyntaxOff
